% =====================================================================
%  第二章  需求分析与总体思路
% =====================================================================
\chapter{需求分析与总体思路}

\section{需求来源与任务拆解}
本课程设计的需求主要来源于两个材料：一是课程给出的 PPT
《课程设计（UNIX 离散化）》，它从页表系统、存储空间管理、离散化存储三个角度
系统阐述了改造目标与修改要求；二是随源码发布的 \texttt{readme.txt}，其中以四点
形式明确了"全面离散化"的具体内涵。综合两者，可将本次课设的总体任务拆解为
表~\ref{tab:task-breakdown} 所列的各项。

\begin{table}[H]
  \centering
  \caption{课设任务拆解}
  \label{tab:task-breakdown}
  \begin{tabular}{clc}
    \toprule
    编号 & 任务项 & 主要涉及文件 \\
    \midrule
    T1 & 修改页表系统，每进程独立页目录 & Process.h / ProcessManager.cpp \\
    T2 & 内存分配改为位示图            & MapNode.h / Allocator.cpp / PageManager.cpp \\
    T3 & 逻辑页离散存放                 & MemoryDescriptor.cpp / PageManager.cpp \\
    T4 & fork 写时复制 (COW)            & ProcessManager.cpp / Exception.cpp \\
    T5 & exec 挂页表 + 参数栈           & ProcessManager.cpp :: Exec \\
    T6 & 正文段共享 (Text 存页框号)     & Text.h / Text.cpp / ProcessManager.cpp \\
    \bottomrule
  \end{tabular}
\end{table}

可以看到，T1\,$\sim$\,T2 属于基础设施改造（页表与分配器），T3\,$\sim$\,T6 则是在
新基础设施之上重构进程生命周期中的内存管理行为。前者是后者的前提，后者是前者
价值的体现。

\section{原系统约束与问题分析}
在着手改造之前，需要厘清原系统在内存管理上存在的约束与问题，它们正是本次改造
所要克服的对象。

\textbf{第一，连续分配带来的外部碎片。}原系统以 \texttt{MapNode} 空闲分区表管理
内存，分配时必须找到一块\emph{地址连续}且\emph{大小足够}的空闲区域。随着进程的
反复创建与退出，物理内存会被切割成大量不连续的小分区，即便总空闲量充足，也可能
因找不到一块足够大的连续区而分配失败。

\textbf{第二，fork 的整体复制浪费。}原系统 \texttt{fork} 创建子进程时，需要把
父进程的整段用户态图像（代码、数据、栈）原样复制一份。而在绝大多数情况下，子
进程要么很快调用 \texttt{exec} 抛弃这份图像，要么只修改其中很小一部分——对整段
图像的立即复制是显著的浪费。

\textbf{第三，盘交换区的复杂度。}原系统依赖 \texttt{SwapperManager} 在内存不足时
把进程图像换出到盘交换区。这一机制需要维护交换空间布局、处理换入换出时序，并且
天然地与"连续进程图像"绑定。本次课设要求\emph{不使用盘交换区}，意味着必须用
请求分页式的按需分配来替代"整体换出"作为内存压力的主要应对手段。

\textbf{第四，相对映射表的局限。}原系统采用一张相对虚实地址映射表作为进程的
物理页表，所有进程共享同一套页表结构，进程切换时需要重写页表项。这一方式难以
支持每个进程独立的地址空间，也与"每进程独立页目录"的目标相悖。

\section{功能需求分析}
针对上述问题，结合 \texttt{readme.txt} 的四点要求，提炼出如下功能需求：

\begin{itemize}[label=\textbf{FR\arabic*}]
  \item \textbf{FR1（逻辑页离散存放）}：任意一个用户虚页都应能映射到任意一个空闲
        的物理页框，不必要求物理页框连续。这是分页机制的核心，也是其余需求的基础。
  \item \textbf{FR2（写时复制 COW）}：\texttt{fork} 后父子进程应共享同一份物理页，
        并将相关页表项置为只读；仅当某一方真正写入时，才为该方分配新物理页并复制
        内容。共享状态需由引用计数维护。
  \item \textbf{FR3（exec 直接挂页表）}：\texttt{exec} 加载新程序时，应为新用户栈
        分配物理页框，在其上接纳命令行参数，然后直接建立用户页表映射，省去把
        整段图像搬运到用户区的内存复制操作。
  \item \textbf{FR4（正文段共享）}：当多个进程执行同一个可执行文件时，应共享同
        一份正文段（代码）物理页。为此，\texttt{Text} 结构中需开数组，记录分配给
        正文段各页的物理页框号。
\end{itemize}

\section{非功能需求与工程约束}
除功能需求外，本设计还受到若干非功能需求与工程约束：

\begin{itemize}
  \item \textbf{不使用盘交换区}：核心约束。所有运行时内存管理路径必须建立在"内存
        中按需分配页框"之上，不允许以换出进程图像来腾挪空间。
  \item \textbf{分配失败不处理}：课设明确不考虑用户区内存分配不成功的情况，上层
        路径可假设分配总能成功。这简化了错误处理，但也在内存不足处理上留下了
        改进空间（见第九章）。
  \item \textbf{保持用户态抽象}：用户程序所见到的虚地址空间布局（text/data/heap/
        stack）应保持不变，离散化对用户程序透明。
  \item \textbf{工程可构建性}：改造后的内核必须能够在本地工具链下完整构建并运行，
        而非停留在代码层面。这要求同时处理工具链差异、镜像构建等问题。
\end{itemize}

\section{总体设计方案}
综合以上分析，本设计的总体方案在数据结构与控制流两个层面展开。

\textbf{数据结构层面}，引入或改造如下核心结构：
\begin{itemize}
  \item \texttt{PageDirectory}：每进程独立页目录，进程创建时分配，记录本进程的
        核心页表、0\# 共享用户页表与 1\# 私有用户页表的物理基址。
  \item 私有用户页表（\texttt{m\_UserPageTableArray}，仅一页）：登记本进程用户态
        各虚页到物理页框的映射。
  \item \texttt{BitMap}：位示图，每个 bit 对应一个 4KB 物理页框，配合
        \texttt{BitMapAllocator} 完成分配与释放。
  \item \texttt{Page[]}：页引用计数数组，记录每个用户区物理页框的共享份数，是
        COW 与 text 共享的物理页生命周期依据。
\end{itemize}

\textbf{控制流层面}，需改造进程生命周期中所有与分页机制相关的操作：
\begin{itemize}
  \item \texttt{NewProc}（fork）：建立子进程私有页目录与私有用户页表，复制父进程
        页表并将共享页置只读、引用计数加一（COW 建立）。
  \item \texttt{PageFault}（缺页异常）：识别写只读共享页的情形，按引用计数决定是
        直接放开写权限还是复制新页（COW 处理）。
  \item \texttt{Exec}：释放旧图像，为新栈分配物理页框并构造参数栈，加载新 text/
        data/stack，建立用户页表映射。
  \item \texttt{Swtch}（进程切换）：通过 \texttt{SwtchUStruct} 更新内核页表 user
        项并刷新 \texttt{CR3} 到目标进程页目录。
  \item \texttt{Exit}：配合引用计数释放数据/栈页，按 \texttt{x\_ccount}/
        \texttt{x\_count} 释放共享 text 段。
\end{itemize}

\section{设计边界与取舍}
为保证改造工作聚焦且可在有限时间内完成，本设计在若干处作出了明确的取舍：

\begin{enumerate}
  \item \textbf{保留 swap 相关遗留路径}。核心用户页框分配、fork/COW、exec/exit 等
        运行时主路径已完全转向离散页框模型，不再依赖连续图像换入换出；但原框架中
        \texttt{SwapperManager} 的若干遗留路径（如 text 副本、僵尸进程 u 区保存）
        并未完全删除，因为它们在当前验证范围内并非主路径。这些遗留路径的进一步
        清理留作后续改进。
  \item \textbf{ATA 采用同步 PIO 替代 DMA}。在离散页表模型下，缓冲区的线性地址
        不能再通过简单掩码直接得到 DMA 所需的物理地址；为保障本地 QEMU 验证的
        可靠性，本设计以同步 PIO 替代 DMA，牺牲一定效率换取正确性。DMA 的真正
        解决方案留待后续。
  \item \textbf{内存不足处理较粗糙}。bitmap 分配器能够判断页框是否可用，但很多
        上层路径仍倾向于假设分配成功，或在失败时直接 panic。更优雅的失败处理与
        回滚策略不在本次课设范围内。
\end{enumerate}

上述取舍在第九章的设计评价中将再次提及。它们划定了本次课设的边界，也指明了后续
可继续完善的方向。
